{{دیوان
| عنوان = تضمین بر شعرِ ابوطالب کلیمؔ
| ماخذ = ویکی ماخذ
| ماخذ_لنک = https://ur.wikisource.org/wiki/%D8%AA%D8%B6%D9%85%DB%8C%D9%86_%D8%A8%D8%B1_%D8%B4%D8%B9%D8%B1%D9%90_%D8%A7%D8%A8%D9%88%D8%B7%D8%A7%D9%84%D8%A8_%DA%A9%D9%84%DB%8C%D9%85
| مقابلہ = اقبال سائبر لائبریری (اقبال اکادمی پاکستان)
| مقابلہ_لنک = https://www.iqbalcyberlibrary.net
}}
<poem>
خوب ہے تجھ کو شعارِ صاحبِ یثرِبؐ کا پاس
کہہ رہی ہے زندگی تیری کہ تُو مسلم نہیں

جس سے تیرے حلقۂ خاتم میں گردُوں تھا اسیر
اے سلیماں! تیری غفلت نے گنوایا وہ نگیں

وہ نشانِ سجدہ جو روشن تھا کوکب کی طرح
ہوگئی ہے اُس سے اب ناآشنا تیری جبیں

دیکھ تو اپنا عمل، تجھ کو نظر آتی ہے کیا
وہ صداقت جس کی بے‌باکی تھی حیرت آفریں

تیرے آبا کی نِگہ بجلی تھی جس کے واسطے
ہے وہی باطل ترے کاشانۂ دل میں مکیں

غافل! اپنے آشیاں کو آ کے پھر آباد کر
نغمہ زن ہے طُورِ معنی پر کلیمِ نکتہ بیں

“سرکشی باہر کہ کردی رامِ او باید شدن
شعلہ ساں از ہر کجا برخاستی، آنجانشیں"
</poem>
